% Standalone integration fragment, based on panel_summary.json and immutable
% recorded source/decode snapshots. Does not claim natural hidden-label truth.
\subsection{Natural precipitation accumulations: a zero-query audit}
\label{sec:noaa-natural}
The aggregate-only information structure also occurs without deliberately
removing fine-resolution labels. NOAA's COOP Hourly Precipitation Data retain
legacy records in which an accumulation total is reported but its distribution
over the intervening hours is unavailable \cite{noaa_hpd_v2,noaa_hpd_format}.
This supports a limited retrospective verification task: determine whether two
fixed forecasts can be ordered for every temporal allocation consistent with
the archived total. It does not establish that historical hourly readings can
be recovered, or assign a cost to recovering them.

We fixed a small source-based panel before inspecting forecast comparisons.
A lexical screen of the first three CONUS station files found nine simple
one-day accumulation records after 2000 at Alberta, Alabama
(USC00010140; $32.2322^{\circ}$N, $87.4104^{\circ}$W). All nine occur in December
2008 or June 2009 and are retained. We also retain every eligible complete day
in those two months. Of 61 calendar dates, 52 are included: nine accumulation
days and 43 complete days. Seven dates have no hourly row, and two contain
multi-day accumulations outside this fixed one-day design. Selection is
purposive and does not estimate the prevalence of useful coarse records.

On each natural case, the hourly record has a beginning flag at 00:00,
interior accumulation flags, and an ending flag with the total at 23:00.
The current CSV definition assigns these fields to $[00{:}00,01{:}00)$ and
$[23{:}00,24{:}00)$ local standard time, respectively. Thus the constraint
covers the full 24-hour recorded day. Values in hundredths of an inch are
converted to millimetres by multiplication by $0.254$. We use the original
accumulation total, never the computed \texttt{DlySum} field as an independent
measurement. The station inventory specifies UTC$-6$ in both months.

The permitted public-source audit includes the parallel 15-minute station
file and the digital raw-data files for both months. For the six December
cases, the parallel file supplies three initial zero quarter-hours followed
by the same unresolved accumulation and matching total. The three June
event dates are absent. No selected station record is present in either
monthly digital raw file. These findings are preserved with source-row
pointers and hashes. Initial zero quarter-hours do not identify any complete
three-hour target: every three-hour bin still contains unconstrained
nonnegative periods, so projecting the checked information onto the eight
bins leaves $y\geq0$ and $\sum_i y_i=T$. Several legacy download endpoints
returned HTTP 404; the conclusion concerns the documented files checked,
not a proof that no other archive exists.

We compare the GEFSv12 control member with the fixed arithmetic mean of its
five daily members, with no fitting or member selection
\cite{guan2022gefs,noaa_gefs_archive}. These are retrospective numerical
weather forecasts, not forecasts issued in 2008--2009 or an ML-performance
benchmark. The nearest $0.25^{\circ}$ grid point is 8.6565\,km from the station.
The nominal cycle is 00:00 UTC; Phase 2 integration actually starts at 03:00
UTC following the replay analysis. Guan et al.'s documented corrections
concern the initial 00:00--03:00 and 00:00--06:00 accumulated fields. We use
valid windows 06:00--06:00 the following day, checked against GRIB validity
metadata, corresponding to eight local three-hour bins. Nested three- and
six-hour precipitation fields are differenced. There were 28 negative
differences, between $-0.05$ and $-0.01$\,mm,
across 27 of 260 member-days. Each discrepancy passed the recorded packing
resolution bound; this establishes consistency with differential quantization,
not its exclusive causal explanation. Our specified nonnegative preprocessing caps the first
three-hour amount at its six-hour total and uses the residual for the second,
preserving every decoded six-hour total. All raw differences and the complete
raw-versus-processed endpoint sensitivity are retained.

Let $a$ denote the processed control and $b$ the processed ensemble mean.
For the daily mean squared-error contrast $\Delta=L(a)-L(b)$, the direct
paired bounds are compared with subtraction of separately sharp individual
MSE bounds using exactly the same information. Table~\ref{tab:noaa-natural}
reports all nine natural cases. Two days certify lower ensemble-mean MSE;
the separate individual bounds certify neither. One day is an exact forecast
tie and six permit both orderings. Uniform allocation gives point orderings
on all eight non-tie days, but six of those orderings are not implied by the
observed information. There is no fine-resolution truth with which to label
those point decisions correct or incorrect.

\begin{table}[t]
\centering
\caption{All natural accumulation cases. Bounds are for control minus
ensemble-mean MSE over eight three-hour bins, in $\mathrm{mm}^2$.
Positive intervals certify the ensemble mean.}
\label{tab:noaa-natural}
\begin{tabular}{lrrrl}
\hline
Date & Total (mm) & Lower & Upper & Conclusion\\
\hline
2008-12-05 & 12.700 & 0.00000 & 0.00000 & exact tie\\
2008-12-10 & 12.700 & 7.18081 & 25.37356 & ensemble\\
2008-12-11 & 33.020 & $-24.65680$ & 26.47467 & unresolved\\
2008-12-17 & 5.080 & $-0.48739$ & 0.43971 & unresolved\\
2008-12-28 & 5.080 & $-1.01836$ & 0.37356 & unresolved\\
2008-12-29 & 5.080 & $-0.04186$ & 0.03942 & unresolved\\
2009-06-05 & 7.874 & $-1.15768$ & 1.48011 & unresolved\\
2009-06-14 & 1.524 & 0.64919 & 1.47063 & ensemble\\
2009-06-24 & 9.652 & $-0.00005$ & 0.04821 & unresolved\\
\hline
\end{tabular}
\end{table}

The conclusions on all nine natural days are unchanged using raw signed
increments; preprocessing changes two dry-control conclusions to ties.
All 43 complete controls have zero recorded precipitation, predominantly
archive-assigned zero flags. Their feasible sets are singleton zero vectors,
so they test only this degenerate case and do not validate nonzero natural
endpoints against hidden truth. Combining their observed loss contrasts with
the nine natural intervals yields the 52-day interval
$[-0.250924,1.195845]\,\mathrm{mm}^2$, which does not order the forecasts.
We also fixed a supplementary $\pm2.54$\,mm total perturbation before inspecting
paired outcomes, motivated by NOAA's nominal $0.1$-inch gauge accuracy.
With $T\in[\max(0,T_0-2.54),T_0+2.54]$, only the 2008-12-10 strict certificate
remains. This is a source-inspired sensitivity, not a validated error bound
for the legacy totals or a confidence statement about physical rainfall.

This panel supplies a concrete natural use of zero-query partial-label
evaluation: some forecast comparisons are identifiable even when their
individual scores are not. Its scope is limited by the nine selected events,
recording and measurement uncertainty, and grid-point versus gauge
representativeness. It provides neither evidence of retrieval savings nor
a general ranking of the forecast systems. The lightweight archive contains
source excerpts, every decoded-message record and provenance pointer, both
forecast variants, and an offline executable reconstruction; fresh GRIB
decoding and full-source checks are distinguished from excerpt arithmetic.
